\paragraph{A smooth proper counterexample.}
\status{Proved and independently audited.}
The failure persists for a smooth proper scheme admitting a smooth proper algebraic lift, and hence persists modulo torsion in test (b).

\begin{proposition}\label{prop:DE-old-proper}
Let $k=\overline{\mathbb F}_3$, let $W=W(k)$, and let $\mathcal E/W$ be the projective smooth elliptic curve
\[
                         y^2z=x^3-xz^2.
\]
Write $E=\mathcal E_k$.  For the literal old site one has
\[
 H^1(\mathcal C_E,\OO)=0,
 \qquad H^1_{\dR}(\mathcal E/W)\simeq W^2.
\]
In particular test (b) is false, including its qualification ``at least modulo torsion.''
\end{proposition}
\begin{proof}
The discriminant of the displayed Weierstrass equation is $64$, a $3$-adic unit.  The invariant differential on its special fibre is a nonzero scalar multiple of $dx/y$.  In characteristic $3$,
\[
 \frac{dx}{y}=\left(\frac{x}{y}\right)^3dx
                  -\left(\frac1y\right)^3x\,dx.
\]
The Cartier operator kills both $dx$ and $x\,dx$, since it selects powers whose exponent is $2$ modulo $3$.  Thus it kills $H^0(E,\Omega^1)$, and, dually, Frobenius on $H^1(E,\OO_E)$ is zero.  The Artin--Schreier sequence, with surjectivity of $a\mapsto a^3-a$ on the algebraically closed field $k$, gives
$H^1_{\et}(E,\mathbb Z/3\mathbb Z)=0$.  Induction using the coefficient exact sequences gives
\[
 H^1_{\et}(E,\mathbb Z/3^n\mathbb Z)=0\qquad(n\geq1).
\]
The finite-level assertion \cite[Corollary~4.7]{Ogus1975}, whose smooth ambient embedding here is $E\hookrightarrow\mathcal E_{W_n}$, therefore gives
\begin{equation}\label{eq:DE-old-proper-infzero}
                  H^1((E/W_n)_{\rm inf},\OO)=0
                  \qquad(n\geq1).
\end{equation}

We explain carefully why this finite-level vanishing detects torsors on the old weak-formal site.  Let $\mathcal T$ be an additive $\OO$-torsor on $\mathcal C_E$, and restrict it to the weak formal completion $\mathcal E^\dagger$.  This means gluing its restrictions on the affine opens, which are old-site objects; the resulting Zariski $\OO$-torsor corresponds to a rank-two extension
\[
     0\longrightarrow\OO_{\mathcal E^\dagger}
       \longrightarrow V\longrightarrow\OO_{\mathcal E^\dagger}
       \longrightarrow0.
\]
For every $N$, the $N$th algebraic diagonal neighbourhood $P_N$ in $\mathcal E\times_W\mathcal E$ is $W$-flat.  Indeed its successive quotients for the diagonal-ideal filtration are symmetric powers of the locally free $\Omega^1_{\mathcal E/W}$.  The two maps from $P_N^\dagger$ to $\mathcal E^\dagger$ give compatible identifications of the two pulled-back extensions, satisfying the cocycle rule on the corresponding triple neighbourhoods.  Thus $V$ carries an ordinary stratification on all these diagonal neighbourhoods.

Reducing modulo $3^n$ turns weak completion into ordinary reduction and gives a stratified extension $V_n$ on the smooth $W_n$-scheme $\mathcal E_n$.  It defines an infinitesimal crystal on $(E/W_n)_{\rm inf}$.  To spell out this standard smooth-ambient construction: for any nilpotent thickening $U\hookrightarrow T$ in that site, formal smoothness locally supplies a map $T\to\mathcal E_n$ lifting $U\to E$.  Pull back $V_n$.  Two choices form a map into a sufficiently high diagonal neighbourhood, so the stratification identifies the pullbacks; the triple cocycle rule gives effective gluing and independence of choices.  The extension and the associated additive torsor are preserved.  Equation~\eqref{eq:DE-old-proper-infzero} says that this stratified extension is split for every $n$.

In particular the underlying class of $\mathcal T|_{\mathcal E^\dagger}$ vanishes modulo every $3^n$.  Meredith's proper comparison theorem, \cite[\S5, Theorem~A, p.~25; or its projective Theorem~B]{Meredith1972}, identifies
\[
 H^1(\mathcal E^\dagger,\OO)
                  \simeq H^1(\mathcal E,\OO)\simeq W.
\]
This module is $3$-adically separated, so the underlying torsor is trivial.  Fix one global trivialization.  The stratification now has an additive cocycle
\[
                         u_N\in H^0(P_N^\dagger,\OO).
\]
Modulo $3^n$, a horizontal trivialization exists by the preceding infinitesimal vanishing.  Its difference from our fixed trivialization is a global function on $\mathcal E_n$, hence an element of $W_n$.  Such a constant has zero diagonal coboundary.  It follows that $u_N$ is zero modulo $3^n$ for every $n$.  Meredith's same proper comparison makes $H^0(P_N^\dagger,\OO)$ a finite $W$-module; $P_N$ is projective and flat.  Hence it is separated, and $u_N=0$ for every $N$.

Finally these diagonal evaluations detect the cocycle on arbitrary old objects, exactly as in Theorem~\ref{thm:DE-old-Gm}; this final step must be included.  Locally choose a smooth affine lift $A$ of an open of $E$ and two maps $a_0,a_1:A\to B$ inducing the same special-fibre map.  Here and below $A$ denotes its weak completion.  Let $J$ be the diagonal ideal and take
\[
                         P_N=(A\otimes_W^\dagger A)/J^N
\]
where $\otimes_W^\dagger$ denotes the weakly completed tensor product; equivalently, weakly complete the algebraic diagonal neighbourhood of the chosen smooth lift.  This is finite over the first $A$ and $W$-flat.  For each $m$, a sufficiently large $N$ makes the pair $(a_0,a_1)$ factor through $P_N\to B/3^mB$.  Indeed the image of $J$ is nilpotent modulo $3$.  The fibre product
\[
                         B\times_{B/3^mB}P_N
\]
is an allowed old-site object by precisely the finiteness, flatness, and square-zero-kernel argument used above.  It forces the cocycle value in $B$ to lie in $3^mB$ for every $m$, hence to vanish.  All local lifts needed here exist: smooth schemes admit local \etale coordinates, one first lifts through the nilpotent kernel modulo $3$, and Hensel's lemma for weakly complete algebras lifts the resulting simple roots; see \cite[Proposition~2.31 and the preceding very-smooth lifting discussion]{DavisLangerZink2012}.  The same argument on common affine opens proves independence of the chosen smooth chart.  Our trivialization on $\mathcal E^\dagger$ therefore descends to a global section of the old-site torsor.  This proves $H^1(\mathcal C_E,\OO)=0$.

Finally, the Hodge exact sequence of the elliptic scheme has two free rank-one outer terms,
\[
 0\longrightarrow H^0(\mathcal E,\Omega^1_{\mathcal E/W})
  \longrightarrow H^1_{\dR}(\mathcal E/W)
  \longrightarrow H^1(\mathcal E,\OO_{\mathcal E})\longrightarrow0.
\]
Thus its middle term is free of rank two.  This proves the second assertion and the claimed failure after killing torsion.
\end{proof}
